\label{sec:literature}

The O-series builds on three distinct strands of the redistricting literature.
This section reviews each strand's key findings and identifies the gaps that the
O-series addresses.

\subsection{Caughey, Tausanovitch, and State Legislative Effects}

\citet{caugheyTausanovitch2017} provide perhaps the most rigorous causal estimate
of redistricting effects on policy outcomes. Using a regression discontinuity design
that exploits close elections at the district margin, they show that the partisanship
of a state legislature---determined in part by its district map---has significant causal
effects on policy liberalism across a broad range of policy areas (taxes, spending,
social policy). A one-standard-deviation increase in Democratic seat share produces
a policy liberalism shift of roughly 0.4 standard deviations on their composite
policy index.

For the O-series, the Caughey-Tausanovitch finding provides the motivation for
the legislative polarization analysis (O.3): if partisan seat share affects policy,
and redistricting affects partisan seat share, then redistricting affects policy
through an ideological selection mechanism. The O.3 paper asks whether this mechanism
operates through polarization specifically---whether compact plans produce legislators
closer to the center, not just different partisan distributions.

The limitation of Caughey and Tausanovitch for the O-series is that their analysis
is at the state legislative, not congressional, level. Congressional redistricting
follows different legal constraints (one-person-one-vote with near-zero tolerance
vs.\ the $\pm5\%$ tolerance permitted in state legislative maps). Congressional
districts are also larger, reducing the intra-district political heterogeneity
that drives the competitive mechanisms in O.1 and O.2.

\subsection{Rodden and Geographic Sorting}

\citet{rodden2019} documents a structural feature of American political geography
that is directly relevant to redistricting outcomes: Democratic voters are systematically
more geographically concentrated (in cities and dense suburbs) than Republican voters.
This geographic sorting creates a structural efficiency disadvantage for Democrats under
any single-member district system: Democratic voters ``waste'' more votes in landslide
districts than Republicans, and this waste is amplified by partisan gerrymandering.

The Rodden framework is important for the O-series in two respects. First, it clarifies
the distinction between geography-induced inefficiency and gerrymandering-induced
inefficiency. Any compact algorithmic plan will inherit some of the Democratic
efficiency disadvantage simply because it must cut between geographically sorted
communities. The O.1 analysis must therefore control for the geographic sorting
baseline when measuring competitiveness effects. Second, Rodden's analysis implies
that the constituent-representative distance effect (O.4) has a partisan asymmetry:
Democratic constituents in urban districts will be closer to their representative
regardless of gerrymandering, because urban representatives locate their offices
in the city center that both the representative and constituents share.

\citet{chen2013} extend the Rodden framework with a simulation-based approach,
generating thousands of random redistricting plans for specific states and showing
that partisan outcomes consistent with Republican advantage arise under most random
plans in states like Florida and North Carolina---not from deliberate gerrymandering
but from geographic sorting. The O-series must control for this baseline to isolate
the gerrymandering contribution to democratic outcomes.

\subsection{Jacobson and the Decline of Competitive Elections}

\citet{jacobson2015} documents the secular decline of competitive House elections
from 1946 to 2014. The fraction of House seats decided by margins under 10 percentage
points declined from roughly 40\% in the 1970s to under 20\% by 2012. Jacobson
attributes this decline to multiple factors: geographic sorting (consistent with Rodden),
incumbency advantage, partisan polarization that reduces ticket-splitting, and
redistricting.

For the O.1 analysis, Jacobson's secular trend presents a confounding challenge:
any redistricting effect on competitiveness must be separated from the secular decline.
We address this by including state-specific time trends in the DiD specification and
by focusing the analysis on redistricting-induced \textit{changes} in competitiveness
rather than levels. States where the 2020 redistricting cycle produced maps with
substantially higher Polsby-Popper scores than the 2010 cycle (naturally or algorithmically)
are the treatment group; states with stable geometry are the control group.

\subsection{Additional Literature}

\textbf{Turnout and competitiveness}: \citet{cox2008} provide a comprehensive review
of the mobilization effects of electoral competition, finding consistent positive effects
across electoral systems. \citet{abramowitz2010} document the decline of swing voters
and its implications for competitive-election mobilization effects. The O.2 analysis
uses these as theoretical priors for the expected sign of the competitiveness-turnout
relationship.

\textbf{Legislative polarization}: \citet{mccarty2006} provide the foundational
DW-NOMINATE based analysis of congressional polarization. \citet{fiorina2008}
debate the extent to which legislative polarization reflects mass public polarization
vs.\ elite selection. The O.3 analysis contributes to this debate by providing
a redistricting-based causal estimate that separates the map contribution from
geographic sorting.

\textbf{Constituent access}: \citet{eckels2020} and \citet{meier1995} provide
the public administration literature on constituency service and district office
location. The O.4 analysis is, to our knowledge, the first systematic empirical
study linking district compactness to constituent travel distance at national scale.

\begin{table}[ht]
\centering
\caption{Prior Literature Summary: O-Series Context}
\label{tab:literature}
\begin{tabular}{llp{6cm}}
\toprule
Study & Outcome & Key Finding \\
\midrule
Caughey \& Tausanovitch (2017) & Policy liberalism & Partisan seat share $\to$ policy (state leg.) \\
Rodden (2019) & Partisan efficiency & Geographic sorting $\to$ Democratic inefficiency \\
Chen \& Rodden (2013) & Competitiveness & Most random plans replicate geographic bias \\
Jacobson (2015) & Competitiveness & Secular decline independent of gerrymandering \\
McCarty et al.\ (2006) & Polarization & DW-NOMINATE gap rising since 1970s \\
Cox \& Munger (1989) & Turnout & Closeness increases turnout via mobilization \\
\bottomrule
\end{tabular}
\end{table}
